\documentclass{article}
\usepackage[utf8]{inputenc}
\usepackage[russian]{babel} % For Russian language support

\title{Пример научной статьи}
\author{Автор}
\date{\today}

\begin{document}

\maketitle

\section{Введение}
Настоящая статья представляет собой пример документа, созданного в формате LaTeX. Введение служит для ознакомления читателя с темой исследования, целями и задачами, а также общей структурой работы. В этой части обычно описываются актуальность темы, основные проблемы, которые рассматриваются в статье, и кратко излагаются ключевые моменты содержания.

LaTeX является мощной системой компьютерной вёрстки, разработанной для подготовки научных и математических документов. Его использование позволяет создавать качественные публикации с правильно оформленными формулами, таблицами и библиографией.

\section{Основная часть}
В основной части статьи рассматриваются ключевые аспекты темы. В зависимости от содержания, здесь могут быть представлены теоретические выкладки, экспериментальные данные, анализ существующих подходов и другие элементы, зависящие от специфики темы исследования.

Для примера, можно рассмотреть основные возможности LaTeX, такие как создание формул:
\begin{equation}
E = mc^2
\end{equation}

А также включение таблиц и списков:
\begin{itemize}
\item Первый пункт
\item Второй пункт
\item Третий пункт
\end{itemize}

\section{Заключение}
В заключении подводятся итоги исследования, формулируются основные выводы и результаты. Здесь также могут быть указаны перспективы дальнейшего изучения темы, возможные направления развития исследуемой области, а также практическое применение полученных результатов.

LaTeX остается важным инструментом для подготовки научных публикаций благодаря своей точности, гибкости и возможностям форматирования сложных документов.

\end{document}